\subsection{Secondary-Objective Ablation}
To isolate the effect of the secondary objectives, both controllers were evaluated offline, through their kinematic models, on the 20 control targets of the multi-section arm (Fig.~S8, Table~S13). For the model-based controller, the control law of Eq.~10 was replayed on the PCC model from the same initial state and target sequence, with and without the null-space term. For the data-driven controller, $N_l^3$ was retrained with $\zeta = 0.01$ (as in Section~3.3.2) and with $\zeta = 0$, with the same data, split, and seed, and evaluated through the frozen $N_p^3$. In both cases the secondary objective reduces $H(\mathbf{l})$ and the deviation of the tendon lengths from $\bar{\mathbf{l}}$. The null-space projection leaves the final position error essentially unchanged, while the penalty in the training loss of $N_l^3$ slightly increases the position error predicted by $N_p^3$. These results are simulated through the respective models and were not measured on the robot.
